\chapter{Computing Limits}

In the last two chapters you learned what a limit \emph{is}. You met the
sharp-edged definition --- for every challenge $\eps$ there is a
response $\delta$ --- and you saw that this machinery exists to pin down a
single idea: \emph{the value a function approaches as the input approaches a
point, regardless of what happens at the point itself.} That was the
foundation. It was also slow. Nobody wants to run an epsilon-delta argument
every time they need a limit, any more than you want to prove the
distributive law every time you multiply two numbers.

This chapter is about \emph{fluency}. We will build a small toolbox ---
substitution, factoring, rationalizing, the squeeze theorem, one-sided limits,
absolute values --- that lets you compute almost any limit you will meet,
quickly and correctly. But this is not a bag of tricks to memorize. Every
tool here is \emph{licensed} by something we already proved: the limit laws.
Each technique is just a clever, legal rearrangement that turns a hard-looking
limit into one you can read off by inspection. The whole art is this:

\begin{focusbox}
Computing a limit means \textbf{transforming the expression --- using only
moves the limit laws permit --- until you reach a form where the answer is
obvious.} The skill is not in the final step. It is in seeing which legal move
clears the obstacle.
\end{focusbox}

We start by recalling the laws, because everything else stands on them.

\section{The Limit Laws: Our License to Compute}
\label{ch07:laws}

Suppose $\displaystyle\lim_{x\to a} f(x) = L$ and
$\displaystyle\lim_{x\to a} g(x) = M$, where $L$ and $M$ are finite real
numbers. In the epsilon-delta chapter we proved the following facts, which
drive everything that follows.

\begin{theorem}[Limit Laws]
\label{ch07:thm-laws}
With the hypotheses above, and $c$ a constant:
\begin{enumerate}
  \item \textbf{(Constant)} $\displaystyle\lim_{x\to a} c = c$.
  \item \textbf{(Identity)} $\displaystyle\lim_{x\to a} x = a$.
  \item \textbf{(Sum/Difference)} $\displaystyle\lim_{x\to a}\big(f(x)\pm g(x)\big) = L \pm M$.
  \item \textbf{(Constant multiple)} $\displaystyle\lim_{x\to a} c\,f(x) = c\,L$.
  \item \textbf{(Product)} $\displaystyle\lim_{x\to a} f(x)\,g(x) = L\,M$.
  \item \textbf{(Quotient)} $\displaystyle\lim_{x\to a} \frac{f(x)}{g(x)} = \frac{L}{M}$, \emph{provided $M \neq 0$}.
  \item \textbf{(Power)} $\displaystyle\lim_{x\to a} \big(f(x)\big)^{n} = L^{n}$ for any positive integer $n$, and $\displaystyle\lim_{x\to a}\sqrt[n]{f(x)} = \sqrt[n]{L}$ when the root is defined.
\end{enumerate}
\end{theorem}

Read these slowly. They say something striking: \emph{the limit of a
combination is the combination of the limits.} Addition, multiplication,
division --- the limit passes straight through them. This is exactly why
calculus is computable at all. If limits did not respect arithmetic, every
single problem would be a fresh epsilon-delta battle.

\begin{whybox}
Why are we allowed to ``pass the limit inside''? The intuition is short. If
$f(x)$ is squeezing in on $L$ and $g(x)$ is squeezing in on $M$, then their
sum $f(x)+g(x)$ has no choice but to squeeze in on $L+M$: push $f$ within
half your tolerance of $L$ and $g$ within half your tolerance of $M$, and the
sum lands within your full tolerance of $L+M$. That ``split the tolerance in
two'' move is the whole proof of the sum law, and the product and quotient
laws are the same spirit with a little more bookkeeping. The laws are not
axioms we hope are true; they are theorems we earned with $\eps$ and $\delta$.
Having earned them, we may now use them freely and never look back.
\end{whybox}

There is one boundary on this freedom, and it is the single most important
caveat in the chapter. \textbf{The quotient law requires $M\neq 0$.} The
moment a denominator's limit is zero, the law goes silent --- it tells you
nothing. That silence is exactly where all the interesting work of this
chapter lives. Most of our techniques are tricks for handling a $0/0$ that
the quotient law refuses to touch.

\begin{trapbox}
A limit law that does not apply is not a law that gives the answer $0$, or
$\infty$, or ``undefined.'' It simply \emph{does not apply}. When you hit
$\frac{0}{0}$ you have not found an answer --- you have found a question. The
expression is called an \emph{indeterminate form} precisely because the laws
determine nothing. You must transform it first.
\end{trapbox}

\section{Direct Substitution: The First Thing to Try}
\label{ch07:substitution}

Here is the happiest fact in limit computation. For the functions you build
out of the everyday pieces --- polynomials, roots, sines, cosines,
exponentials, logs, and any sums, products, quotients, and compositions of
these --- the limit as $x\to a$ is just the value at $a$, \emph{whenever that
value makes sense}. You plug in. You are done.

\begin{theorem}[Substitution for continuous functions]
\label{ch07:thm-subst}
If $f$ is continuous at $a$, then
\[
  \lim_{x\to a} f(x) = f(a).
\]
Every polynomial is continuous everywhere. Rational functions, roots,
$\exp$, $\log$, $\sin$, $\cos$, and their compositions are continuous
everywhere they are defined.
\end{theorem}

We will give continuity its own careful chapter. For now take the working
definition: \emph{$f$ is continuous at $a$ when the limit there equals the
value there}, with no gap, no jump, and no hole. The theorem then looks almost
circular --- ``the limit is the value when the limit is the value'' --- but
the content is in the second sentence: \emph{the standard building blocks
really are continuous}, and that fact is itself a consequence of the limit
laws.

\begin{whybox}
Why does substitution work for a polynomial? Watch the laws assemble it. Take
$p(x) = 3x^{2} - 5x + 1$ and $a = 2$. By the identity law $\lim x = 2$. By the
power law $\lim x^{2} = 2^{2} = 4$. By the constant-multiple law
$\lim 3x^{2} = 12$ and $\lim 5x = 10$. By the sum/difference law the whole
thing is $12 - 10 + 1 = 3$, which is exactly $p(2)$. Nothing here was a leap
of faith --- each step cited a law from
Theorem~\ref{ch07:thm-laws}. Substitution is not a shortcut that bypasses the
laws; it is the laws, applied so routinely we stopped narrating them.
\end{whybox}

\begin{example}[Substitution, straight ahead]
\label{ch07:ex-subst1}
\[
  \lim_{x\to 3} \frac{x^{2}+1}{x-1}
  = \frac{3^{2}+1}{3-1}
  = \frac{10}{2}
  = 5.
\]
The denominator's limit is $2\neq 0$, so the quotient law applies and we
simply evaluate. No drama.
\end{example}

\begin{example}[Substitution through a composition]
\label{ch07:ex-subst2}
\[
  \lim_{x\to 0} \cos\!\big(e^{x} - 1\big)
  = \cos\!\big(e^{0} - 1\big)
  = \cos(0) = 1.
\]
The inside $e^{x}-1 \to 0$, and $\cos$ is continuous at $0$, so we evaluate
through the composition.
\end{example}

\begin{trapbox}
Substitution is the \emph{first} thing to try, not the only thing. If plugging
in gives a genuine number, you are finished. If it gives $\frac{0}{0}$,
$\frac{\text{nonzero}}{0}$, $\frac{\infty}{\infty}$, or any other
nonsense, \textbf{stop} --- you have learned that the function is not
continuous at $a$ (or not even defined there), and you must reach for a
different tool. Substitution failing is not failure; it is a diagnosis.
\end{trapbox}

\begin{figure}[htbp]
\centering
\begin{tikzpicture}
\begin{axis}[
  width=10cm, height=6.2cm,
  title style={font=\small},
  title={Direct substitution: continuous curves pass through $(a, f(a))$},
  axis lines=middle,
  xlabel={$x$}, ylabel={$f(x)$},
  xmin=1.0, xmax=4.2, ymin=-0.5, ymax=7,
  xtick={3}, ytick={5},
  xticklabels={$3$}, yticklabels={$5$},
  enlargelimits=false,
  samples=120, domain=1.6:4.1,
]
\addplot[curveblue, very thick] {(x^2+1)/(x-1)};
\draw[densely dashed, gray] (axis cs:3,0) -- (axis cs:3,5);
\draw[densely dashed, gray] (axis cs:1.0,5) -- (axis cs:3,5);
\addplot[only marks, mark=*, tangentred] coordinates {(3,5)};
\end{axis}
\end{tikzpicture}
\caption{For $\dfrac{x^{2}+1}{x-1}$ near $x=3$ the curve is smooth and
unbroken. The limit is simply the height of the curve there: $5$. This is the
easy case --- substitution and nothing more.}
\label{ch07:fig-subst}
\end{figure}

\section{When Substitution Fails: The 0/0 Form}
\label{ch07:zerozero}

The interesting limits --- the ones calculus actually needs --- almost always
arrive as $\frac{0}{0}$, for a structural reason. The derivative itself,
which we are marching toward, is
\[
  \lim_{h\to 0} \frac{f(x+h) - f(x)}{h},
\]
and as $h\to 0$ the top goes to $f(x)-f(x)=0$ while the bottom goes to $0$.
\emph{Every derivative is a $0/0$ limit.} So learning to handle $0/0$ is not a
sideshow; it is the central engine of differential calculus, rehearsed here
on simpler functions.

\begin{happeningbox}
What is really happening in a $\frac{0}{0}$ limit? Both the numerator and the
denominator are shrinking to zero, and the limit measures \emph{how fast each
shrinks compared to the other}. If the top heads to zero twice as fast as the
bottom, the ratio approaches $2$. If they vanish at exactly the same
rate, you get a nonzero number; if the top wins, you get $0$; if the bottom
wins, you get a blow-up. The symbol $\frac{0}{0}$ hides all of this --- it is
``indeterminate'' precisely because the bare zeros throw away the rate
information. Our algebra tricks are nothing but ways to \emph{recover the
hidden rate} by canceling the common factor that is causing both to vanish.
\end{happeningbox}

The key structural fact is this. When a rational expression gives $\frac00$ at
$x=a$, it is because $(x-a)$ is a factor of \emph{both} the numerator and the
denominator. That shared factor is the culprit: it is what drives both to zero.
Cancel it, and the limit becomes visible.

\begin{whybox}
Why is canceling a factor legitimate when the factor is zero in the limit?
Here is the subtle, beautiful point. The functions
\[
  \frac{x^2 - 9}{x - 3} \qquad\text{and}\qquad x + 3
\]
are \emph{not the same function}: the first is undefined at $x=3$, the second
is perfectly happy there. But a limit as $x\to 3$ \textbf{never looks at
$x=3$ itself} --- only at the values \emph{near} it. For every $x\neq 3$ the
two expressions are identical, because dividing $x^2-9=(x-3)(x+3)$ by $(x-3)$
is legal whenever $x-3\neq 0$. Since the limit ignores the single point $x=3$,
the two functions have the same limit there. The cancellation is licensed not
by the quotient law (which is silent) but by the very \emph{definition} of a
limit: limits care about the neighborhood, not the point. This is the deepest
``why'' in the chapter, so reread it until it feels obvious.
\end{whybox}

\begin{figure}[htbp]
\centering
\begin{tikzpicture}
\begin{axis}[
  width=10cm, height=6cm,
  title style={font=\small},
  title={Canceling the zero-factor: exposing the continuous line},
  axis lines=middle,
  xlabel={$x$}, ylabel={$y$},
  xmin=-1, xmax=5, ymin=-1, ymax=9,
  xtick={3}, ytick={6}, xticklabels={$3$}, yticklabels={$6$},
  samples=2, domain=-0.8:4.8,
]
\addplot[curveblue, very thick] {x+3};
\draw[densely dashed, gray] (axis cs:3,0) -- (axis cs:3,6);
\draw[densely dashed, gray] (axis cs:0,6) -- (axis cs:3,6);
\addplot[only marks, mark=*, fill=white, draw=curveblue, line width=1pt] coordinates {(3,6)};
\node[anchor=west, tangentred] at (axis cs:3.15,6.6) {hole at $x=3$};
\end{axis}
\end{tikzpicture}
\caption{The graph of $\dfrac{x^{2}-9}{x-3}$ is the straight line $y=x+3$ with
a single \emph{puncture} at $x=3$. The function is undefined there, but the
limit reads the height the line is heading for: $6$. The missing point is
invisible to the limit.}
\label{ch07:fig-hole}
\end{figure}

\subsection{Technique 1: Factor and Cancel}
\label{ch07:factor}

The recipe writes itself: when substitution gives $\frac00$ in a rational
expression, factor top and bottom, cancel the common $(x-a)$, then substitute
into what remains.

\begin{example}[The canonical cancellation]
\label{ch07:ex-factor1}
\[
  \lim_{x\to 3} \frac{x^{2}-9}{x-3}.
\]
Substitution gives $\frac{0}{0}$. Factor the numerator:
\[
  \frac{x^{2}-9}{x-3} = \frac{(x-3)(x+3)}{x-3} = x+3
  \quad\text{for } x\neq 3.
\]
Now substitution into $x+3$ is legal: $\displaystyle\lim_{x\to 3}(x+3) = 6.$
\end{example}

\begin{example}[Factoring a quadratic over a quadratic]
\label{ch07:ex-factor2}
\[
  \lim_{x\to 2} \frac{x^{2}-x-2}{x^{2}-4}.
\]
At $x=2$ both top and bottom vanish, so $(x-2)$ divides both. Factor:
\[
  \frac{x^{2}-x-2}{x^{2}-4}
  = \frac{(x-2)(x+1)}{(x-2)(x+2)}
  = \frac{x+1}{x+2}
  \quad (x\neq 2).
\]
Substitute: $\displaystyle \lim_{x\to 2}\frac{x+1}{x+2} = \frac{3}{4}.$
\end{example}

\begin{example}[A cube; remember the factoring patterns]
\label{ch07:ex-factor3}
\[
  \lim_{x\to 1}\frac{x^{3}-1}{x-1}.
\]
Use $x^{3}-1 = (x-1)(x^{2}+x+1)$:
\[
  \frac{x^{3}-1}{x-1} = x^{2}+x+1 \quad (x\neq 1),
  \qquad
  \lim_{x\to 1}(x^{2}+x+1) = 3.
\]
\end{example}

\begin{labbox}
Try this one before reading on:
\[
  \lim_{x\to -2}\frac{x^{2}+5x+6}{x+2}.
\]
Substitution gives $\frac{0}{0}$, so factor: $x^{2}+5x+6=(x+2)(x+3)$. Cancel
the $(x+2)$ to get $x+3$, then substitute $x=-2$ to find the limit is $1$. Did
you also notice that the answer is the \emph{slope-like} number you'd get from
the rate idea? The numerator and denominator both vanish at $x=-2$, and the
ratio of their vanishing rates is exactly $1$.
\end{labbox}

\subsection{Technique 2: Rationalize}
\label{ch07:rationalize}

Square roots block factoring. The cure is the \emph{conjugate}. Recall the
difference-of-squares identity $(\sqrt{A}-\sqrt{B})(\sqrt{A}+\sqrt{B}) = A-B$:
multiplying a root expression by its conjugate turns the roots into ordinary
subtraction, which usually exposes the hidden $(x-a)$ factor.

\begin{whybox}
Why does multiplying by the conjugate help? A subtraction of roots,
$\sqrt{x+1}-2$, is opaque --- you cannot see what factor is making it vanish.
But $(\sqrt{x+1}-2)(\sqrt{x+1}+2) = (x+1)-4 = x-3$ is transparent: the root is
gone and the vanishing factor $(x-3)$ stares back at you. We multiply top and
bottom by the same conjugate, so we have multiplied by $1$ and changed
nothing --- but we have \emph{relocated} the difficulty from the unreadable
numerator to a readable one. It is the conjugate's job to convert
``difference of roots'' into ``difference of the things under the roots.''
\end{whybox}

\begin{example}[Rationalizing the numerator]
\label{ch07:ex-rat1}
\[
  \lim_{x\to 3}\frac{\sqrt{x+1}-2}{x-3}.
\]
Substitution gives $\frac{0}{0}$ (the top is $\sqrt4-2=0$). Multiply top and
bottom by the conjugate $\sqrt{x+1}+2$:
\[
  \frac{\sqrt{x+1}-2}{x-3}\cdot\frac{\sqrt{x+1}+2}{\sqrt{x+1}+2}
  = \frac{(x+1)-4}{(x-3)\big(\sqrt{x+1}+2\big)}
  = \frac{x-3}{(x-3)\big(\sqrt{x+1}+2\big)}.
\]
Now the $(x-3)$ cancels:
\[
  = \frac{1}{\sqrt{x+1}+2}
  \quad (x\neq 3),
  \qquad
  \lim_{x\to 3}\frac{1}{\sqrt{x+1}+2} = \frac{1}{\sqrt4+2} = \frac14.
\]
\end{example}

\begin{example}[Conjugate on the denominator]
\label{ch07:ex-rat2}
\[
  \lim_{x\to 0}\frac{x}{\sqrt{x+4}-2}.
\]
Multiply by the conjugate $\sqrt{x+4}+2$ over itself:
\[
  \frac{x}{\sqrt{x+4}-2}\cdot\frac{\sqrt{x+4}+2}{\sqrt{x+4}+2}
  = \frac{x\big(\sqrt{x+4}+2\big)}{(x+4)-4}
  = \frac{x\big(\sqrt{x+4}+2\big)}{x}
  = \sqrt{x+4}+2 \quad (x\neq 0).
\]
Then $\displaystyle\lim_{x\to 0}\big(\sqrt{x+4}+2\big) = 2+2 = 4.$
\end{example}

\subsection{Technique 3: Clear Complex Fractions}
\label{ch07:complexfrac}

A fraction buried inside a fraction looks fierce but yields to the oldest
trick in algebra: combine the little fractions over a common denominator, then
simplify. These appear constantly in derivative computations.

\begin{example}[A fraction of fractions]
\label{ch07:ex-complex}
\[
  \lim_{x\to 2}\frac{\dfrac{1}{x}-\dfrac{1}{2}}{x-2}.
\]
Combine the numerator over the common denominator $2x$:
\[
  \frac{1}{x}-\frac{1}{2} = \frac{2-x}{2x}.
\]
So the whole expression is
\[
  \frac{2-x}{2x}\cdot\frac{1}{x-2}
  = \frac{-(x-2)}{2x(x-2)}
  = \frac{-1}{2x} \quad (x\neq 2),
\]
and $\displaystyle\lim_{x\to 2}\frac{-1}{2x} = \frac{-1}{4}.$ Notice the sign
flip from $2-x = -(x-2)$ --- a classic place to lose a minus.
\end{example}

\begin{trapbox}
The recurring danger in all three techniques is \textbf{sign errors} during
factoring, especially $a-b = -(b-a)$. When you cancel, watch which way the
factor points. A single dropped minus sign turns a correct answer into its
negative, and the algebra will look perfectly clean while it lies to you.
\end{trapbox}

\section{The Squeeze Theorem}
\label{ch07:squeeze}

Factoring and rationalizing handle algebraic obstructions. But some limits ---
the ones involving wild oscillation, or transcendental functions like $\sin$
near $0$ --- resist algebra entirely. For these we need a different idea, and
it happens to be one of the most intuitive in analysis.

\begin{intuitionbox}
Imagine a small fish swimming between two larger fish, one above and one below.
If the top fish and the bottom fish both swim toward the same doorway and meet
there, the little fish, trapped between them, has no choice --- it must go
through that doorway too. The squeeze theorem is exactly this. If a function
is trapped between two others that converge to a common limit, it is forced to
that same limit, no matter how erratically it wiggles in between.
\end{intuitionbox}

\begin{theorem}[Squeeze Theorem]
\label{ch07:thm-squeeze}
Suppose that for all $x$ near $a$ (except possibly at $a$ itself)
\[
  g(x) \le f(x) \le h(x),
\]
and that
\[
  \lim_{x\to a} g(x) = \lim_{x\to a} h(x) = L.
\]
Then $\displaystyle\lim_{x\to a} f(x) = L$ as well.
\end{theorem}

\begin{whybox}
Why is the squeeze theorem true? An epsilon-delta argument makes it airtight,
but the idea is the picture. Fix a tolerance $\eps$. Since $g\to L$, eventually
$g(x) > L-\eps$. Since $h\to L$, eventually $h(x) < L+\eps$. Once both of these
hold (take the smaller of the two $\delta$'s), the sandwich $g\le f\le h$ forces
\[
  L-\eps < g(x) \le f(x) \le h(x) < L+\eps,
\]
so $f(x)$ is within $\eps$ of $L$. That is precisely the definition of
$\lim f = L$. The bread does the work; $f$ is just along for the ride.
\end{whybox}

\begin{example}[Taming an oscillation]
\label{ch07:ex-squeeze-osc}
\[
  \lim_{x\to 0} x^{2}\sin\!\frac{1}{x}.
\]
Substitution is hopeless: $\sin(1/x)$ oscillates infinitely fast as $x\to 0$
and has no limit. But $\sin$ of anything lives in $[-1,1]$, so
\[
  -1 \le \sin\frac1x \le 1
  \quad\Longrightarrow\quad
  -x^{2} \le x^{2}\sin\frac1x \le x^{2}.
\]
Both ends, $-x^2$ and $x^2$, go to $0$ as $x\to 0$. Squeezed between them, so
does our function:
\[
  \lim_{x\to 0} x^{2}\sin\frac1x = 0.
\]
The $x^{2}$ factor is a clamp that crushes the wild oscillation down to zero.
\end{example}

\begin{figure}[htbp]
\centering
\begin{tikzpicture}
\begin{axis}[
  width=11cm, height=6cm,
  title style={font=\small},
  title={Squeeze theorem: bounding curves clamp wild oscillations},
  axis lines=middle,
  xlabel={$x$}, ylabel={$y$},
  xmin=-0.55, xmax=0.55, ymin=-0.32, ymax=0.32,
  samples=200, domain=-0.5:0.5,
  legend pos=north west, legend style={font=\footnotesize},
]
\addplot[gray, dashed, thick] {x^2};
\addlegendentry{$y=x^2$}
\addplot[gray, dashed, thick] {-x^2};
\addlegendentry{$y=-x^2$}
\addplot[curveblue, very thick, samples=200, domain=-0.5:-0.004] {x^2*sin(deg(1/x))};
\addplot[curveblue, very thick, samples=200, domain=0.004:0.5] {x^2*sin(deg(1/x))};
\addplot[only marks, mark=*, fill=white, draw=tangentred, line width=1pt] coordinates {(0,0)};
\end{axis}
\end{tikzpicture}
\caption{$y = x^{2}\sin(1/x)$ oscillates faster and faster near $0$, but it is
pinned between the parabolas $y=\pm x^{2}$ (dashed). Since both clamps close on
$0$, the trapped curve is forced to $0$.}
\label{ch07:fig-squeeze}
\end{figure}

\subsection{The Crown Jewel: $\sin x / x \to 1$}
\label{ch07:sinx}

We now use the squeeze theorem to prove the most important single limit in all
of calculus:
\[
  \boxed{\;\lim_{x\to 0}\frac{\sin x}{x} = 1.\;}
\]
This is the limit from which \emph{every} derivative of a trig function flows.
When you later learn that $\ddx{\sin x} = \cos x$, the engine inside that fact
is this one limit. It is worth proving carefully.

\begin{intuitionbox}
For a small angle $x$ (in radians!), the arc of the unit circle and the
straight chord beneath it are nearly the same length, and $\sin x$ is almost
the arc itself. ``Nearly'' becomes ``exactly in the limit.'' The claim
$\sin x/x \to 1$ is the precise statement that, for tiny angles, \emph{the sine
of an angle is essentially the angle}. This is why radians matter: only in
radians does arc length equal the angle, making the comparison clean.
\end{intuitionbox}

\begin{figure}[htbp]
\centering
\begin{tikzpicture}[scale=3.3]
  \coordinate (O) at (0,0);
  \coordinate (A) at (1,0);
  \coordinate (B) at (38:1);
  \coordinate (F) at ({cos(38)},0);
  \coordinate (T) at (1,{tan(38)});
  \draw[->] (-0.15,0) -- (1.25,0);
  \draw[->] (0,-0.15) -- (0,1.15);
  % arc on unit circle (length = x)
  \draw[curveblue, very thick] (A) arc (0:38:1);
  % radii
  \draw (O) -- (A);
  \draw (O) -- (B);
  \draw (O) -- (T);
  % sin segment (vertical drop from B)
  \draw[targetgreen, very thick] (F) -- (B);
  % tangent segment
  \draw[quotientorange, very thick] (A) -- (T);
  \node at (0.18,0.07) {$x$};
  \node[below] at (A) {$1$};
  \fill (B) circle (0.012);
  \fill (A) circle (0.012);
  \node[targetgreen, left] at ({cos(38)},{sin(38)/2}) {$\sin x$};
  \node[quotientorange, right] at (T) {$\tan x$};
  \node[curveblue, right] at (19:1) {arc $=x$};
\end{tikzpicture}
\caption{The unit circle, angle $x$ in radians. The vertical green segment is
$\sin x$, the bold blue arc has length $x$, and the orange segment up to the
tangent line has length $\tan x$. Comparing the three areas gives
$\sin x \le x \le \tan x$.}
\label{ch07:fig-unitcircle}
\end{figure}

\begin{focusbox}
\textbf{The proof, in one sentence:} we trap $x$ between $\sin x$ and $\tan x$
using three nested areas in the unit circle, divide through to isolate
$\dfrac{\sin x}{x}$, and squeeze.
\end{focusbox}

\textbf{Setting up the inequality.} Work in the unit circle (radius $1$) and
take a small positive angle $x$, with $0 < x < \frac{\pi}{2}$. Look at
Figure~\ref{ch07:fig-unitcircle}. We compare three regions:
\begin{itemize}
  \item the \emph{small triangle} $OAB$ with the chord, of area $\tfrac12\sin x$;
  \item the \emph{circular sector} $OAB$, of area $\tfrac12 x$;
  \item the \emph{large triangle} $OAT$ out to the tangent line, of area $\tfrac12\tan x$.
\end{itemize}
Let us justify each area. The small triangle $OAB$ has base $OA=1$ along the
$x$-axis and height equal to the vertical drop from $B$, which is $\sin x$; so
its area is $\tfrac12\cdot 1\cdot \sin x = \tfrac12\sin x$. The sector is the
fraction $\tfrac{x}{2\pi}$ of the full disk of area $\pi(1)^2 = \pi$, giving
area $\tfrac{x}{2\pi}\cdot\pi = \tfrac12 x$. The large triangle $OAT$ has base
$OA=1$ and height $AT = \tan x$ (the tangent line meets the vertical at height
$\tan x$), so its area is $\tfrac12\tan x$.

Now read off the geometry: the small triangle sits \emph{inside} the sector,
which sits \emph{inside} the large triangle. Hence
\[
  \underbrace{\tfrac12\sin x}_{\text{triangle}}
  \;\le\;
  \underbrace{\tfrac12 x}_{\text{sector}}
  \;\le\;
  \underbrace{\tfrac12\tan x}_{\text{triangle}}.
\]
Multiply through by $2$:
\[
  \sin x \;\le\; x \;\le\; \tan x = \frac{\sin x}{\cos x}.
\]

\textbf{Rearranging into a sandwich.} Since $0<x<\frac\pi2$ we have
$\sin x > 0$, so we may divide the whole chain by $\sin x$ without flipping any
inequalities:
\[
  1 \;\le\; \frac{x}{\sin x} \;\le\; \frac{1}{\cos x}.
\]
Take reciprocals --- which \emph{does} reverse the inequalities, since all
quantities are positive:
\[
  \cos x \;\le\; \frac{\sin x}{x} \;\le\; 1.
\]
There it is: $\dfrac{\sin x}{x}$ is sandwiched between $\cos x$ and $1$
(Figure~\ref{ch07:fig-sinx}).

\textbf{Squeezing.} As $x\to 0^{+}$, the lower bread $\cos x \to \cos 0 = 1$,
and the upper bread is the constant $1$. Both converge to $1$, so by the
squeeze theorem
\[
  \lim_{x\to 0^{+}}\frac{\sin x}{x} = 1.
\]

\textbf{The other side.} We assumed $x>0$. What about $x\to 0^{-}$? The
function $\dfrac{\sin x}{x}$ is \emph{even}: replacing $x$ by $-x$ gives
$\dfrac{\sin(-x)}{-x} = \dfrac{-\sin x}{-x} = \dfrac{\sin x}{x}$, unchanged. So
the left-hand limit equals the right-hand limit. Both are $1$, and therefore
the two-sided limit is
\[
  \lim_{x\to 0}\frac{\sin x}{x} = 1. \qquad\blacksquare
\]

\begin{figure}[htbp]
\centering
\begin{tikzpicture}
\begin{axis}[
  width=11cm, height=6cm,
  title style={font=\small},
  title={The squeeze sandwich: $\cos x \le \frac{\sin x}{x} \le 1$ around $x=0$},
  axis lines=middle,
  xlabel={$x$}, ylabel={$y$},
  xmin=-3.3, xmax=3.3, ymin=-0.35, ymax=1.25,
  ytick={1}, yticklabels={$1$},
  samples=200, domain=-3.2:3.2,
  legend pos=south east, legend style={font=\footnotesize},
]
\addplot[gray, dashed, thick, domain=-1.5:1.5] {cos(deg(x))};
\addlegendentry{$y=\cos x$ (lower bread)}
\addplot[gray, dotted, very thick] coordinates {(-3.2,1) (3.2,1)};
\addlegendentry{$y=1$ (upper bread)}
\addplot[curveblue, very thick, domain=-3.2:-0.03] {sin(deg(x))/x};
\addplot[curveblue, very thick, domain=0.03:3.2] {sin(deg(x))/x};
\addplot[only marks, mark=*, fill=white, draw=tangentred, line width=1pt] coordinates {(0,1)};
\end{axis}
\end{tikzpicture}
\caption{The function $y=\dfrac{\sin x}{x}$ (solid) is pinned between
$y=\cos x$ (dashed, shown only near $0$ where it is the lower bound) and the
line $y=1$ (dotted). As $x\to 0$ both breads close on $1$, forcing the curve up
to the punctured value $1$. The point $x=0$ is a removable hole: the formula is
undefined there, but the limit is plainly $1$.}
\label{ch07:fig-sinx}
\end{figure}

\begin{historybox}
The fact that $\sin x \approx x$ for small angles was known and \emph{used}
long before anyone could prove it with limits. Indian astronomers of the
Kerala school (Madhava, around 1400) worked with series for sine that encode
exactly this small-angle behavior, and Newton and Euler computed fluently with
$\sin x / x \to 1$ in the late 1600s and 1700s while treating the ``$x$'' as a
genuine infinitesimal --- a number smaller than any positive number yet not
zero. Their instinct was flawless and their answers correct, but the
foundation was shaky: what \emph{is} an infinitesimal? It was Cauchy in the
1820s who recast such statements as rigorous limits, and the squeeze theorem ---
the area argument you just saw --- is essentially his way of replacing the
hand-wavy ``$\sin x$ and $x$ become equal'' with a rigorous bound that anyone
can check. The geometry was the ancients'; the rigor is the nineteenth
century's.
\end{historybox}

\begin{example}[A companion limit]
\label{ch07:ex-cos}
A close relative, proved with the same machinery and a trig identity:
\[
  \lim_{x\to 0}\frac{1-\cos x}{x} = 0.
\]
Multiply by the conjugate $1+\cos x$:
\[
  \frac{1-\cos x}{x}\cdot\frac{1+\cos x}{1+\cos x}
  = \frac{1-\cos^{2}x}{x(1+\cos x)}
  = \frac{\sin^{2}x}{x(1+\cos x)}
  = \frac{\sin x}{x}\cdot\frac{\sin x}{1+\cos x}.
\]
As $x\to 0$ the first factor $\to 1$ and the second $\to \frac{0}{2} = 0$, so
the product $\to 1\cdot 0 = 0$. Two of our tricks --- conjugates and the
$\sin x/x$ limit --- working together.
\end{example}

\begin{example}[Using $\sin x/x$ as a building block]
\label{ch07:ex-sin-build}
\[
  \lim_{x\to 0}\frac{\sin 5x}{3x}.
\]
The trick is to manufacture the form $\frac{\sin(\text{something})}{\text{same
something}}$. Multiply and divide to make the argument of sine match its
denominator:
\[
  \frac{\sin 5x}{3x}
  = \frac{\sin 5x}{5x}\cdot\frac{5x}{3x}
  = \frac{\sin 5x}{5x}\cdot\frac{5}{3}.
\]
As $x\to 0$, $5x\to 0$ too, so $\dfrac{\sin 5x}{5x}\to 1$, and the limit is
$1\cdot\dfrac53 = \dfrac53.$
\end{example}

\section{One-Sided Limits}
\label{ch07:onesided}

So far we have let $x$ approach $a$ from both sides at once. But sometimes a
function behaves differently coming from the left than from the right, and we
need to look at each approach separately. These are \emph{one-sided limits}.

\begin{definition}[One-sided limits]
\label{ch07:def-onesided}
The \textbf{right-hand limit} $\displaystyle\lim_{x\to a^{+}} f(x) = L$ means
$f(x)$ approaches $L$ as $x$ approaches $a$ through values \emph{greater} than
$a$. The \textbf{left-hand limit} $\displaystyle\lim_{x\to a^{-}} f(x) = L$
uses values \emph{less} than $a$. The epsilon-delta phrasing is identical to
the two-sided one, except we restrict to $0 < x - a < \delta$ (right) or
$-\delta < x - a < 0$ (left).
\end{definition}

The connection between one-sided and two-sided limits is the whole reason we
care about them, and it is clean:

\begin{theorem}[Two-sided $=$ both one-sided agreeing]
\label{ch07:thm-onesided}
$\displaystyle\lim_{x\to a} f(x) = L$ \emph{if and only if}
\[
  \lim_{x\to a^{-}} f(x) = L \quad\text{and}\quad \lim_{x\to a^{+}} f(x) = L.
\]
That is: \textbf{the two-sided limit exists exactly when both one-sided limits
exist and are equal.}
\end{theorem}

\begin{intuitionbox}
A two-sided limit is a meeting. For the function to ``approach $L$ as $x\to a$''
it must approach $L$ whether you sneak up from the left or from the right. If
the two approaches head for different heights, there is no single number the
function is closing in on --- so the two-sided limit \emph{does not exist}.
This is the standard way to prove a limit fails to exist: show the two sides
disagree.
\end{intuitionbox}

\begin{example}[A jump: the limit fails to exist]
\label{ch07:ex-jump}
The Heaviside step function
\[
  H(x) = \begin{cases} 0 & x < 0,\\ 1 & x \ge 0,\end{cases}
\]
has, at $x=0$,
\[
  \lim_{x\to 0^{-}} H(x) = 0,
  \qquad
  \lim_{x\to 0^{+}} H(x) = 1.
\]
The two sides disagree, so $\displaystyle\lim_{x\to 0} H(x)$ \textbf{does not
exist.} The graph jumps; no single height is being approached.
\end{example}

\begin{figure}[htbp]
\centering
\begin{tikzpicture}
\begin{axis}[
  width=9cm, height=4.6cm,
  title style={font=\small},
  title={Step discontinuity: left limit $0 \ne$ right limit $1$},
  axis lines=middle,
  xlabel={$x$}, ylabel={$H(x)$},
  xmin=-2.2, xmax=2.2, ymin=-0.4, ymax=1.6,
  ytick={1}, xtick={\empty},
]
\addplot[curveblue, very thick] coordinates {(-2.1,0) (-0.02,0)};
\addplot[curveblue, very thick] coordinates {(0,1) (2.1,1)};
\addplot[only marks, mark=*, fill=white, draw=curveblue, line width=1pt] coordinates {(0,0)};
\addplot[only marks, mark=*, fill=curveblue, draw=curveblue] coordinates {(0,1)};
\end{axis}
\end{tikzpicture}
\caption{The Heaviside step. From the left, $H\to 0$; from the right,
$H\to 1$. Because the two one-sided limits disagree, the two-sided limit at
$0$ does not exist.}
\label{ch07:fig-heaviside}
\end{figure}

\begin{example}[Different formulas on each side]
\label{ch07:ex-piecewise}
Let
\[
  f(x) = \begin{cases} x^{2}+1 & x \le 1,\\ 4 - 2x & x > 1.\end{cases}
\]
From the left we use $x^2+1$: $\displaystyle\lim_{x\to 1^{-}} f(x) = 1^2+1 = 2$.
From the right we use $4-2x$: $\displaystyle\lim_{x\to 1^{+}} f(x) = 4-2 = 2$.
The sides agree, so $\displaystyle\lim_{x\to 1} f(x) = 2$ \emph{exists},
regardless of the fact that the formula switches at $x=1$
(Figure~\ref{ch07:fig-piecewise}).
\end{example}

\begin{figure}[htbp]
\centering
\begin{tikzpicture}
\begin{axis}[
  width=9.5cm, height=5.4cm,
  title style={font=\small},
  title={Piecewise agreement: two formulas meeting at the same height},
  axis lines=middle,
  xlabel={$x$}, ylabel={$f(x)$},
  xmin=-1.6, xmax=3.2, ymin=-0.5, ymax=4.2,
  xtick={1}, ytick={2}, xticklabels={$1$}, yticklabels={$2$},
  samples=80,
]
\addplot[curveblue, very thick, domain=-1.4:1] {x^2+1};
\addplot[tangentred, very thick, domain=1:3.1] {4-2*x};
\draw[densely dashed, gray] (axis cs:1,0) -- (axis cs:1,2);
\addplot[only marks, mark=*, fill=curveblue, draw=curveblue] coordinates {(1,2)};
\end{axis}
\end{tikzpicture}
\caption{The piecewise $f$: the parabola $x^2+1$ (blue, for $x\le 1$) and the
line $4-2x$ (red, for $x>1$) meet at the same height $2$ when $x=1$. Both
one-sided limits equal $2$, so the two-sided limit exists even though the
formula changes. Contrast this with the Heaviside jump, where the pieces fail
to meet.}
\label{ch07:fig-piecewise}
\end{figure}

\begin{walkbox}
Take a walk and picture this. A function can have a perfectly good two-sided
limit at a point where it is built from \emph{two different formulas}, as long
as those formulas meet at the same height. And a function can \emph{fail} to
have a limit at a point where it is given by a single tidy formula --- if that
formula blows up or jumps. The limit cares only about the \emph{heights being
approached from each side}, never about how the function is written or what it
does at the point itself. Can you sketch a function continuous everywhere
except one jump, and one that is defined by two formulas yet perfectly smooth?
\end{walkbox}

\section{Limits with Absolute Values}
\label{ch07:absvalue}

The absolute value $\abs{x}$ is secretly a piecewise function in disguise:
\[
  \abs{x} = \begin{cases} x & x \ge 0,\\ -x & x < 0,\end{cases}
\]
and that hidden split is \emph{exactly} the place where one-sided limits earn
their keep. Whenever an absolute value sits on top of (or inside) an
expression, the right move is almost always: \textbf{break it at the point
where its inside changes sign, and take the limit from each side.}

\begin{example}[The classic $\abs{x}/x$]
\label{ch07:ex-absx}
\[
  \lim_{x\to 0}\frac{\abs{x}}{x}.
\]
For $x>0$, $\abs{x}=x$, so $\frac{\abs x}{x} = \frac{x}{x} = 1$. For $x<0$,
$\abs{x}=-x$, so $\frac{\abs x}{x} = \frac{-x}{x} = -1$. Therefore
\[
  \lim_{x\to 0^{+}}\frac{\abs{x}}{x} = 1,
  \qquad
  \lim_{x\to 0^{-}}\frac{\abs{x}}{x} = -1.
\]
The two sides disagree, so the two-sided limit \textbf{does not exist}. The
function is $+1$ on the right and $-1$ on the left --- a jump straight through
the origin (Figure~\ref{ch07:fig-absx}).
\end{example}

\begin{figure}[htbp]
\centering
\begin{tikzpicture}
\begin{axis}[
  width=9cm, height=4.6cm,
  title style={font=\small},
  title={Absolute value jump: $\frac{|x|}{x}$ switches sign at $0$},
  axis lines=middle,
  xlabel={$x$}, ylabel={$\frac{\abs{x}}{x}$},
  xmin=-2.2, xmax=2.2, ymin=-1.6, ymax=1.6,
  ytick={-1,1}, xtick={\empty},
]
\addplot[curveblue, very thick] coordinates {(-2.1,-1) (-0.02,-1)};
\addplot[curveblue, very thick] coordinates {(0.02,1) (2.1,1)};
\addplot[only marks, mark=*, fill=white, draw=curveblue, line width=1pt] coordinates {(0,-1) (0,1)};
\end{axis}
\end{tikzpicture}
\caption{The sign function $\dfrac{\abs{x}}{x}$ equals $+1$ to the right of $0$
and $-1$ to the left, with the point $x=0$ removed (the function is undefined
there). The left and right limits, $-1$ and $+1$, disagree, so the two-sided
limit at $0$ does not exist.}
\label{ch07:fig-absx}
\end{figure}

\begin{happeningbox}
What is the absolute value really doing to a limit? It \emph{folds the negative
side up}, which reverses the direction of approach on the left. So an
expression that is heading one way as $x$ increases through $0$ from the right
may head the opposite way from the left --- the fold flips its sign. That is
why absolute-value limits so often split into two unequal one-sided limits.
The cure is never to fight the absolute value head-on but to \emph{unfold it}
into its two cases and treat each side as the plain function it really is.
\end{happeningbox}

\begin{example}[Absolute value not at zero]
\label{ch07:ex-abs-shift}
\[
  \lim_{x\to 3}\frac{\abs{x-3}}{x-3}.
\]
The inside $x-3$ changes sign at $x=3$. For $x>3$, $x-3>0$ so the ratio is
$+1$; for $x<3$, $x-3<0$ so the ratio is $-1$. Right limit $=1$, left limit
$=-1$: the two-sided limit does not exist. (It is the previous example shifted
three units to the right.)
\end{example}

\begin{example}[When the absolute value does \emph{not} cause trouble]
\label{ch07:ex-abs-ok}
\[
  \lim_{x\to 0}\, x\,\abs{x}.
\]
Here splitting is easy and the sides agree. For $x\ge 0$, $x\abs x = x^2 \to 0$;
for $x<0$, $x\abs x = -x^2 \to 0$. Both one-sided limits are $0$, so the
two-sided limit is $0$. An absolute value does not automatically destroy a
limit --- it only matters when the two folded sides head to different places.
\end{example}

\begin{trapbox}
Do not ``simplify'' $\abs{x}$ to $x$ and barge ahead. $\abs{x} = x$ is true
only for $x\ge 0$. On the left of the trouble point you must use $\abs{x}=-x$,
and forgetting the sign is the single most common error here. When you see an
absolute value in a limit, your reflex should be: \emph{find where the inside
is zero, then split.}
\end{trapbox}

\section{One-Sided Limits and Infinite Blow-Ups}
\label{ch07:infinite}

One-sided limits also let us describe what happens when a function blows up ---
the $\frac{\text{nonzero}}{0}$ situation that the quotient law also refuses to
touch, but for a different reason than $\frac00$.

\begin{example}[Opposite infinities from each side]
\label{ch07:ex-blowup}
\[
  \lim_{x\to 0}\frac{1}{x}.
\]
Substitution gives $\frac{1}{0}$ --- not indeterminate, but undefined.
Approach from the right ($x>0$, tiny): $\frac1x$ is a huge positive number, so
$\displaystyle\lim_{x\to 0^{+}}\frac1x = +\infty$. From the left ($x<0$, tiny):
$\frac1x$ is a huge \emph{negative} number, so
$\displaystyle\lim_{x\to 0^{-}}\frac1x = -\infty$. The two sides shoot off in
opposite directions, so the two-sided limit does not exist.
\end{example}

\begin{example}[Same infinity from both sides]
\label{ch07:ex-blowup2}
\[
  \lim_{x\to 0}\frac{1}{x^{2}} = +\infty.
\]
Here $x^2 > 0$ no matter which side you come from, so $\frac{1}{x^2}\to+\infty$
from \emph{both} sides. We write $+\infty$ to describe the shared behavior,
while remembering that ``$=\infty$'' is shorthand for ``grows without bound,''
not a claim that the limit is a number.
\end{example}

\begin{trapbox}
$\frac{\text{nonzero}}{0}$ and $\frac{0}{0}$ are completely different
diagnoses. A nonzero number over zero \emph{blows up} (vertical asymptote);
there is nothing to cancel and no finite limit. A zero over zero is
\emph{indeterminate}; there is a hidden common factor and the limit may be any
finite number. When substitution fails, your very first question is: \emph{is
the numerator also zero?} The answer routes you to entirely different
techniques.
\end{trapbox}

\section{A Preview of Indeterminate Forms}
\label{ch07:indeterminate}

We have spent this chapter conquering one indeterminate form, $\frac00$, with
algebra. Its siblings recur throughout calculus, and each demands its own
tools, so they are worth naming now. An \emph{indeterminate form} is a
symbolic ``shape'' a limit can take where the limit laws give no answer ---
the value genuinely depends on the competing rates, not on the shape alone.

\begin{focusbox}
The seven classic indeterminate forms:
\[
  \frac{0}{0},\quad \frac{\infty}{\infty},\quad 0\cdot\infty,\quad
  \infty-\infty,\quad 1^{\infty},\quad 0^{0},\quad \infty^{0}.
\]
Each is ``indeterminate'' because the answer can be \emph{anything} depending
on how fast the pieces move. The shape alone does not decide.
\end{focusbox}

Why are these indeterminate while, say, $\frac{2}{0}$ or $0^{5}$ are not? The
difference is whether the symbolic form pins down the answer.

\begin{intuitionbox}
Consider $0\cdot\infty$. One factor is shrinking to nothing, the other growing
without bound --- they are at war, and who wins depends on \emph{how fast} each
moves. Watch three skirmishes as $x\to 0^{+}$, each of the form
$(\text{something}\to 0)\cdot(\text{something}\to\infty)$:
\[
  x\cdot\frac1x \to 1,
  \qquad
  x\cdot\frac{1}{x^{2}} \to +\infty,
  \qquad
  x^{2}\cdot\frac1x \to 0.
\]
Same shape $0\cdot\infty$, three different answers. That is the entire meaning
of ``indeterminate'': the form is a question, never an answer. By contrast
$\frac{2}{0}$ always blows up and $0^{5}$ is always $0$ --- those forms
\emph{are} determined, so we do not call them indeterminate.
\end{intuitionbox}

The cleanest power tool for the two leading forms, $\frac00$ and
$\frac{\infty}{\infty}$, is \emph{l'H\^opital's rule}, which turns the limit of
a ratio into the limit of the ratio of derivatives. It is wonderfully
effective --- and it is worthless until you have derivatives, which is why it
waits for its own chapter much later. The forms $0\cdot\infty$,
$\infty-\infty$, and the three exponential shapes $1^\infty$, $0^0$,
$\infty^0$ are all handled by first \emph{algebraically rewriting} them as
$\frac00$ or $\frac{\infty}{\infty}$ (the exponential ones by taking logs), and
then applying the same machinery. The algebra you practiced in this chapter ---
factoring, rationalizing, clearing fractions --- is the warm-up; l'H\^opital is
the heavy artillery you get once derivatives are in hand.

\begin{connectionbox}
There is a thread running from this chapter all the way to Taylor series near
the end of the book. The reason $\frac{\sin x}{x}\to 1$ is that $\sin x = x -
\frac{x^3}{6} + \cdots$; divide by $x$ and the leading term is $1$ with the
rest vanishing. Every indeterminate $\frac00$ is, underneath, a competition
between the \emph{leading terms} of two series. The algebra tricks of this
chapter, l'H\^opital's rule, and Taylor expansion are three faces of the same
idea: \emph{a limit is decided by the lowest-order behavior of each piece.}
\end{connectionbox}

\section{A Strategy, Start to Finish}
\label{ch07:strategy}

Put the tools in order. Faced with $\displaystyle\lim_{x\to a} f(x)$:

\begin{enumerate}
  \item \textbf{Try substitution.} Plug in $a$. If you get a finite number and
  the function is continuous there, you are done.
  \item \textbf{Diagnose the failure.} If substitution gives
  $\frac{\text{nonzero}}{0}$, it is a blow-up --- use one-sided limits to
  describe $\pm\infty$ and conclude the two-sided limit (usually) does not
  exist. If it gives $\frac00$, it is indeterminate --- proceed.
  \item \textbf{Remove the common factor.} Factor and cancel; if there are
  roots, rationalize with the conjugate; if there are nested fractions, clear
  them. Then substitute into the simplified form.
  \item \textbf{Squeeze if algebra fails.} For oscillations or trig limits like
  $\frac{\sin x}{x}$, trap the function between two others with a common limit.
  \item \textbf{Split at sign changes.} For absolute values or piecewise
  formulas, take left and right limits separately and check whether they agree.
  \item \textbf{Recognize the deeper forms.} If you are left with
  $\frac\infty\infty$, $0\cdot\infty$, or an exponential shape, name the
  indeterminate form and remember l'H\^opital is coming.
\end{enumerate}

\begin{figure}[htbp]
\centering
\begin{tikzpicture}[
  level 1/.style={sibling distance=4.5cm, level distance=1.3cm},
  node/.style={draw=curveblue, fill=curveblue!8, rounded corners=3pt, font=\footnotesize, align=center, inner sep=4pt, line width=0.7pt},
  leaf/.style={draw=#1, fill=#1!10, rounded corners=3pt, font=\scriptsize, align=center, inner sep=3.5pt, line width=0.7pt},
  edge from parent/.style={draw=gray!70, -Stealth, thick}
]
\node[node, font=\small\bfseries] {Evaluate $\lim_{x\to a} f(x)$ by substitution $f(a)$}
  child {
    node[leaf=targetgreen] {\textbf{Finite Real $L$}\\[2pt]Continuous at $a$\\$\implies \lim = f(a)$}
  }
  child {
    node[leaf=quotientorange] {\textbf{Form $c/0$ ($c\neq 0$)}\\[2pt]Vertical asymptote\\Test $x\to a^\pm$ for $\pm\infty$}
  }
  child {
    node[leaf=tangentred] (indet) {\textbf{Form $0/0$ Indeterminate}\\[2pt]Algebraic transformation}
  };

% Child actions for 0/0
\node[leaf=curveblue, below left=0.8cm and 0.4cm of indet] (fact) {\textbf{Polynomials:}\\Factor \& cancel $(x-a)$};
\node[leaf=structurepurple, below=0.8cm of indet] (conj) {\textbf{Radicals:}\\Multiply by conjugate};
\node[leaf=tangentred, below right=0.8cm and 0.4cm of indet] (trig) {\textbf{Trig / $|x|$ / Squeeze:}\\$\frac{\sin x}{x}\to 1$ or split $a^\pm$};

\draw[edge from parent] (indet) -- (fact);
\draw[edge from parent] (indet) -- (conj);
\draw[edge from parent] (indet) -- (trig);
\end{tikzpicture}
\caption{The limit computation decision roadmap: diagnose the initial form upon substitution, then deploy the matching algebraic or geometric tool to resolve the obstruction.}
\label{ch07:fig-decision-tree}
\end{figure}

\begin{example}[A limit that uses several tools]
\label{ch07:ex-combo}
\[
  \lim_{x\to 0}\frac{\sqrt{1+x}-1}{\sin x}.
\]
Substitution gives $\frac00$. There is a root, so rationalize the numerator,
and there is a sine, so plan to use $\frac{\sin x}{x}$. Multiply the numerator
by its conjugate:
\[
  \frac{\sqrt{1+x}-1}{\sin x}\cdot\frac{\sqrt{1+x}+1}{\sqrt{1+x}+1}
  = \frac{(1+x)-1}{\sin x\,(\sqrt{1+x}+1)}
  = \frac{x}{\sin x\,(\sqrt{1+x}+1)}.
\]
Now split the factors deliberately:
\[
  = \frac{x}{\sin x}\cdot\frac{1}{\sqrt{1+x}+1}
  = \frac{1}{\frac{\sin x}{x}}\cdot\frac{1}{\sqrt{1+x}+1}.
\]
As $x\to 0$, $\frac{\sin x}{x}\to 1$ so its reciprocal $\to 1$, and
$\sqrt{1+x}+1 \to 2$. The limit is
\[
  1\cdot\frac{1}{2} = \frac12.
\]
Rationalizing and the $\sin x/x$ limit, working side by side.
\end{example}

\begin{aestheticbox}
Notice the shape of what you have learned. Every technique in
this chapter does the same thing: it \emph{changes the appearance of an
expression without changing the value it approaches}, until the answer becomes
self-evident. Factoring, conjugates, common denominators, the squeeze
sandwich, splitting at a sign change --- these are not five unrelated tricks
but five dialects of a single sentence: \emph{a limit lives in the
neighborhood of a point, not at the point, so I am free to rewrite the function
on that neighborhood however I please.} That freedom, licensed by the limit
laws and the definition of a limit itself, is the quiet engine of all of
calculus. Once you feel it, computing limits stops being a chore and becomes
a search for the right disguise to remove.
\end{aestheticbox}

\section*{Exercises}
\addcontentsline{toc}{section}{Exercises}

\begin{exercise}[Substitution]
\label{ch07:exr-subst}
Evaluate by direct substitution, justifying that it is legal:
\[
  \text{(a)}\ \lim_{x\to 2}\big(3x^{2}-4x+7\big),
  \qquad
  \text{(b)}\ \lim_{x\to 0}\frac{e^{x}+\cos x}{x+2},
  \qquad
  \text{(c)}\ \lim_{x\to \pi}\sin\!\Big(\frac{x}{2}\Big).
\]
\end{exercise}

\begin{exercise}[Factor and cancel]
\label{ch07:exr-factor}
Evaluate:
\[
  \text{(a)}\ \lim_{x\to 4}\frac{x^{2}-16}{x-4},
  \qquad
  \text{(b)}\ \lim_{x\to -1}\frac{x^{2}+3x+2}{x^{2}-1},
  \qquad
  \text{(c)}\ \lim_{x\to 2}\frac{x^{3}-8}{x-2}.
\]
\end{exercise}

\begin{exercise}[Rationalize]
\label{ch07:exr-rat}
Evaluate:
\[
  \text{(a)}\ \lim_{x\to 0}\frac{\sqrt{9+x}-3}{x},
  \qquad
  \text{(b)}\ \lim_{x\to 1}\frac{\sqrt{x}-1}{x-1},
  \qquad
  \text{(c)}\ \lim_{x\to 7}\frac{\sqrt{x+2}-3}{x-7}.
\]
\end{exercise}

\begin{exercise}[Complex fractions]
\label{ch07:exr-complex}
Evaluate:
\[
  \text{(a)}\ \lim_{x\to 3}\frac{\frac{1}{x}-\frac13}{x-3},
  \qquad
  \text{(b)}\ \lim_{h\to 0}\frac{\frac{1}{(2+h)^{2}}-\frac14}{h}.
\]
\end{exercise}

\begin{exercise}[Trig limits]
\label{ch07:exr-trig}
Using $\displaystyle\lim_{x\to 0}\frac{\sin x}{x}=1$ and
$\displaystyle\lim_{x\to 0}\frac{1-\cos x}{x}=0$, evaluate:
\[
  \text{(a)}\ \lim_{x\to 0}\frac{\sin 7x}{2x},
  \qquad
  \text{(b)}\ \lim_{x\to 0}\frac{\sin 3x}{\sin 5x},
  \qquad
  \text{(c)}\ \lim_{x\to 0}\frac{\tan x}{x}.
\]
\end{exercise}

\begin{exercise}[Squeeze]
\label{ch07:exr-squeeze}
(a) Show $\displaystyle\lim_{x\to 0} x^{4}\cos\!\frac{1}{x^{2}} = 0$ by
squeezing. (b) Given that $\displaystyle 1 - \frac{x^{2}}{6} \le
\frac{\sin x}{x} \le 1$ for small $x$, use the squeeze theorem to re-derive
$\displaystyle\lim_{x\to 0}\frac{\sin x}{x}=1$.
\end{exercise}

\begin{exercise}[One-sided and absolute value]
\label{ch07:exr-onesided}
Evaluate each one-sided limit and decide whether the two-sided limit exists:
\[
  \text{(a)}\ \lim_{x\to 5}\frac{\abs{x-5}}{x-5},
  \qquad
  \text{(b)}\ \lim_{x\to 2^{+}}\frac{x-2}{\abs{x-2}},
  \qquad
  \text{(c)}\ \lim_{x\to 0}\frac{x^{2}}{\abs{x}}.
\]
\end{exercise}

\begin{exercise}[Piecewise]
\label{ch07:exr-piecewise}
Let
\[
  f(x) = \begin{cases} 2x+1 & x < 1,\\ k & x = 1,\\ x^{2}+2 & x > 1.\end{cases}
\]
Find $\displaystyle\lim_{x\to 1^{-}}f(x)$ and
$\displaystyle\lim_{x\to 1^{+}}f(x)$. Does $\displaystyle\lim_{x\to 1}f(x)$
exist? Does the value of $k$ affect your answer? Explain in one sentence why
the limit ignores $k$.
\end{exercise}

\begin{exercise}[Blow-ups]
\label{ch07:exr-blowup}
Describe the behavior with one-sided limits, and state whether the two-sided
limit exists:
\[
  \text{(a)}\ \lim_{x\to 2}\frac{1}{x-2},
  \qquad
  \text{(b)}\ \lim_{x\to 2}\frac{1}{(x-2)^{2}},
  \qquad
  \text{(c)}\ \lim_{x\to 1}\frac{x}{x^{2}-1}.
\]
\end{exercise}

\begin{exercise}[Conceptual]
\label{ch07:exr-concept}
Decide whether each statement is true or false, and justify in a sentence.
(a) If $\lim_{x\to a} f(x)$ exists then $f$ is defined at $a$.
(b) If $f(a)$ is defined then $\lim_{x\to a} f(x) = f(a)$.
(c) The form $\frac00$ always has limit $1$ because ``anything over itself is
$1$.''
(d) If both one-sided limits exist, the two-sided limit exists.
\end{exercise}

\begin{exercise}[Proof: a squeeze]
\label{ch07:exr-proof1}
Prove using the squeeze theorem that if $\abs{f(x)} \le M\abs{x-a}$ for all $x$
near $a$ (with $M$ a fixed constant), then $\displaystyle\lim_{x\to a} f(x)=0$.
\emph{Hint:} sandwich $f$ between $-M\abs{x-a}$ and $M\abs{x-a}$.
\end{exercise}

\begin{exercise}[Proof: the conjugate, in general]
\label{ch07:exr-proof2}
Let $a>0$. Prove that
\[
  \lim_{x\to a}\frac{\sqrt{x}-\sqrt{a}}{x-a} = \frac{1}{2\sqrt{a}}
\]
by rationalizing. Explain in one sentence which limit law licenses the final
substitution step, and why it is now legal when it was not before.
\end{exercise}
